\documentclass[11pt, a4paper]{article}
\usepackage{kotex}
\usepackage{geometry}
\usepackage{booktabs}
\usepackage{amsmath}
\usepackage{listings}
\usepackage{xcolor}
\usepackage{hyperref}
\usepackage{enumitem}
\usepackage{fancyhdr}
\usepackage{mdframed}

\geometry{margin=2.5cm}

\hypersetup{colorlinks=true, linkcolor=blue!60!black}

\definecolor{codegreen}{rgb}{0,0.5,0}
\definecolor{codegray}{rgb}{0.4,0.4,0.4}
\definecolor{codepurple}{rgb}{0.5,0,0.5}
\definecolor{backcolour}{rgb}{0.97,0.97,0.97}
\definecolor{keywordblue}{rgb}{0,0,0.7}

\lstdefinestyle{matlab}{
    backgroundcolor=\color{backcolour},
    commentstyle=\color{codegreen}\itshape,
    keywordstyle=\color{keywordblue}\bfseries,
    numberstyle=\tiny\color{codegray},
    stringstyle=\color{codepurple},
    basicstyle=\ttfamily\small,
    breakatwhitespace=false,
    breaklines=true,
    captionpos=b,
    keepspaces=true,
    numbers=left,
    numbersep=6pt,
    showspaces=false,
    showstringspaces=false,
    showtabs=false,
    tabsize=4,
    language=Matlab,
    frame=single,
    rulecolor=\color{black!20},
    xleftmargin=\parindent,
}

\lstset{style=matlab}

\pagestyle{fancy}
\fancyhf{}
\rhead{PSW Dispersion Technical Report}
\lhead{\leftmark}
\cfoot{\thepage}

\title{\textbf{PSW Dispersion 시뮬레이션}\\
\textbf{리팩토링 및 병렬화 --- 기술 참조 보고서}\\[6pt]
\large 코드 수정 전체 내역 및 설계 근거}
\author{}
\date{2026년 4월}

\begin{document}
\maketitle
\tableofcontents
\newpage

%% ---------------------------------------------------------------------------
\section{개요}

\subsection{작업 목표}

원래 순차 실행 코드(\texttt{Main\_code.m})를 다음 목표로 리팩토링하였다:

\begin{itemize}
  \item 8개 $I_2$ 세기 조건 $\times$ 200,000 분자를 병렬화
  \item 코드를 기능별 함수 파일로 분리하여 재사용성 향상
  \item 불필요한 반복 계산 및 루프 제거
\end{itemize}

\subsection{하드웨어 환경}

\begin{center}
\begin{tabular}{ll}
\toprule
구성 & 사양 \\
\midrule
CPU & Intel Core i7-13700K (P-core 8$\times$2 + E-core 8 = 논리 24코어) \\
RAM & DDR5 64 GB \\
GPU & NVIDIA RTX 4060 Ti (VRAM 8 GB 전용 + 공유 32 GB, CUDA cores 4352) \\
MATLAB & R2026a \\
\bottomrule
\end{tabular}
\end{center}

\subsection{최종 파일 구조}

\begin{verbatim}
PSW dispersion/
  run_all.m                  <- entry point (new)
  STN_constants.mat          <- physical constants (unchanged)
  linearly.mat               <- cos2theta lookup table (unchanged)
  data/
    alpha_T1_J=2n,j,m.txt   <- quantum state data (unchanged)
  sim/
    init_particles.m         <- particle initializer (new)
    build_aligned_cost.m     <- vectorized aligned_cost (new)
    time_integrate.m         <- Velocity Verlet integrator (new)
    axay_opt_partial.m       <- optical force, refactored (new)
    image_processing.m       <- image processing parfor (new)
    voigt2d_aniso_rot.m      <- PSF function (unchanged)
\end{verbatim}

\subsection{성능 결과 요약}

\begin{center}
\begin{tabular}{lrr}
\toprule
방식 & 시간 & 배율 \\
\midrule
기존 순차 (8조건 $\times$ 120 s) & $\sim$960 s & 1$\times$ \\
병렬화 최적화 (8 workers) & 181 s & 5.3$\times$ \\
\bottomrule
\end{tabular}
\end{center}

%% ---------------------------------------------------------------------------
\section{병렬화 전략 설계}

\subsection{Sub-Batch Parfor 방식 선택 근거}

처음에는 GPU(gpuArray) 활용을 검토하였으나 다음 이유로 CPU parfor를 채택하였다:

\begin{mdframed}[backgroundcolor=backcolour, linecolor=black!20]
\textbf{GPU 성능 문제 분석}\\
MATLAB \texttt{gpuArray}는 element-wise 연산마다 별도 CUDA 커널을 launch한다.
\texttt{axay\_opt\_partial} 내 연산 수 $\times$ 시간 적분 6,000 스텝
= 수만 건의 kernel launch overhead가 발생하여 실측 결과 GPU 84 s vs CPU 64 s로
GPU가 더 느렸다. (20,000 입자 기준; 입자 수가 충분히 크면 GPU가 유리해질 수 있음)
\end{mdframed}

\subsection{Sub-Batch 분해 설계}

\[
  \text{총 작업} = \underbrace{200{,}000}_{\text{MAX/조건}} \times \underbrace{8}_{\text{n\_groups}}
  \;\longrightarrow\;
  \underbrace{20{,}000}_{\text{batch\_size}} \times \underbrace{80}_{\text{n\_tasks}}
\]

각 태스크 번호 \texttt{task} $\in [1, 80]$에서 해당 조건 인덱스로의 매핑:
\[
  g = \left\lceil \frac{\text{task}}{n\_\text{sub}} \right\rceil, \quad
  s = \bigl((\text{task}-1) \bmod n\_\text{sub}\bigr) + 1
\]

\subsection{메모리 대역폭 병목 발견}

\texttt{parfor} worker 수를 늘리면서 실측한 결과:

\begin{center}
\begin{tabular}{rrr}
\toprule
Workers & 시간 (s) & 전 단계 대비 \\
\midrule
8  & \textbf{181} & --- \\
12 & 183 & +2 \\
16 & 187 & +4 \\
24 & 208 & +21 \\
\bottomrule
\end{tabular}
\end{center}

Worker를 늘릴수록 오히려 느려지는 이유:
각 worker가 \texttt{aligned\_cost} 배열(20,000 $\times$ nI, broadcast 변수)을
시간 루프 매 스텝 접근하므로, 동시 접근이 DDR5 메모리 버스를 포화시킨다.
\textbf{최적: 8 workers}.

\subsection{MaxNumWorkers 강제 확장 문제}

MATLAB \texttt{parpool}의 local 프로파일은 \texttt{MaxNumWorkers} 기본값이 8이다.
단순히 \texttt{parpool(N)}만 호출하면 8개 이상 연결되지 않는다.
\texttt{parcluster} API를 통해 프로파일을 직접 수정해야 한다:

\begin{lstlisting}
cluster = parcluster('local');
if cluster.NumWorkers < n_workers
    cluster.NumWorkers = n_workers;
    saveProfile(cluster);   % save profile changes
end
parpool(cluster, n_workers);
\end{lstlisting}

%% ---------------------------------------------------------------------------
\section{파일별 수정 상세}

\subsection{\texttt{run\_all.m} (신규 작성)}

기존 \texttt{Main\_code.m}을 완전히 대체하는 진입점.

\subsubsection{주요 변경사항}

\begin{enumerate}[itemsep=4pt]
  \item \textbf{Parallel pool 설정 블록 추가}: 기존 pool이 부족한 경우 재시작 로직 포함
  \item \textbf{MAX 오버라이드}: \texttt{STN\_constants.mat}의 MAX=20,000을
        200,000으로 오버라이드
  \item \textbf{linearly\_2d 사전 로드}: broadcast 변수로 모든 worker에 공유
        (\texttt{parfor} 내부에서 매 task마다 로드하지 않도록)
  \item \textbf{jm\_data 경로 수정}: 상수 파일에는 파일명만 저장되어 있으므로
        \texttt{fullfile(root\_dir, 'data', jm\_data)}로 full path 구성
  \item \textbf{이미지 처리 파라미터 하드코딩}: \texttt{STN\_constants.mat}에 없는
        값들(\texttt{rep\_std}, \texttt{rep\_vel}, \texttt{vel\_pix},
        \texttt{atten\_factorx}, \texttt{atten\_factory})을
        \texttt{STN\_dis\_CS2\_velocity.mlx}에서 추출하여 직접 설정

  \item \textbf{sub-batch parfor 구현}:
\begin{lstlisting}
results = cell(n_tasks, 1);
parfor task = 1:n_tasks
    g_idx = ceil(task / n_sub);
    [particle_t, laser_t] = init_particles(jm_data_path, batch_size, ...);
    j_idx_t = particle_t(:,10) + 1;
    m_idx_t = particle_t(:,11) + 1;
    aligned_cost_t = linearly_2d(sub2ind([nj_lin,nm_lin], ...
                                         j_idx_t, m_idx_t), :);
    I1_t = repmat(I1_val,         batch_size, 1);
    I2_t = repmat(I2_vals(g_idx), batch_size, 1);
    results{task} = time_integrate(..., false);  % false = CPU mode
end
\end{lstlisting}

  \item \textbf{결과 재조립}: task 결과를 조건별로 \texttt{vertcat}으로 합산
\begin{lstlisting}
Cell_velocity = cell(n_groups, 1);
for g = 1:n_groups
    task_idx = (g-1)*n_sub + (1:n_sub);
    Cell_velocity{g} = vertcat(results{task_idx});  % 200,000 x 3
end
\end{lstlisting}
\end{enumerate}

\subsection{\texttt{sim/axay\_opt\_partial.m} (리팩토링)}

기존 \texttt{axay\_opt\_partial.m}에서 다음을 수정하였다.

\subsubsection{수정 1: 오타 수정}

\begin{center}
\begin{tabular}{ll}
\toprule
기존 (오타) & 수정 후 \\
\midrule
\texttt{acceleartion} & \texttt{acceleration} \\
\texttt{alpa}         & \texttt{alpha} \\
\bottomrule
\end{tabular}
\end{center}

\subsubsection{수정 2: gauss\_t 사전 계산 (성능)}

\begin{lstlisting}
% Before: exp(-4*log(2)*t.^2/tau^2) computed 6x redundantly

% Fix: precompute once
gauss_t = exp(-4 * log(2) * t.^2 / tau^2);  % MAX x 1
% reused throughout the function
\end{lstlisting}

\subsubsection{수정 3: inv\_mass 상수 정의}

\begin{lstlisting}
inv_mass = 6.022e23 / 76.14e-3;
% CS2 molar mass 76.14 g/mol, Avogadro's number
\end{lstlisting}

\subsubsection{수정 4: col\_idx 경계 클램핑}

룩업 테이블 열 인덱스가 범위를 벗어나는 경우 방어:

\begin{lstlisting}
n_cols  = size(aligned_value, 2);
col_idx = min(I_floor + 1, n_cols - 1);  % upper clamp
col_idx = max(col_idx, 1);               % lower clamp (prevent 0)
\end{lstlisting}

\subsubsection{수정 5: GPU 호환 allocation}

\begin{lstlisting}
% Before: zeros(MAX, 3)  -> CPU only
% Fix:
acceleration = zeros(MAX, 3, 'like', position);
% if position is gpuArray -> acceleration is also gpuArray
% if position is a regular array -> same behavior
\end{lstlisting}

\subsection{\texttt{sim/build\_aligned\_cost.m} (신규)}

기존 for 루프를 완전히 대체하는 벡터화 함수.

\subsubsection{기존 코드}

\begin{lstlisting}
% O(MAX) for loop - repeated MAX=200,000 times
for index = 1:MAX
    aligned_cost(index,:) = linearly(particle(index,10)+1, ...
                                     particle(index,11)+1, :);
end
\end{lstlisting}

\subsubsection{수정 후}

\begin{lstlisting}
function aligned_cost = build_aligned_cost(particle, linearly, MAX)
    [nj, nm, ~] = size(linearly);
    linearly_2d = reshape(linearly, nj * nm, []);   % (nj*nm) x nI
    j_idx = particle(:,10) + 1;                      % MAX x 1
    m_idx = particle(:,11) + 1;                      % MAX x 1
    lin_idx = sub2ind([nj, nm], j_idx, m_idx);       % MAX x 1
    aligned_cost = linearly_2d(lin_idx, :);           % MAX x nI
end
\end{lstlisting}

단일 벡터화 인덱싱으로 교체하여 수십$\sim$수백 배 속도 향상.
\textbf{주의}: \texttt{run\_all.m}에서는 \texttt{linearly\_2d}를 이미 broadcast 변수로
사전 생성하므로 이 함수 대신 직접 인덱싱을 사용한다.

\subsection{\texttt{sim/init\_particles.m} (신규)}

기존 코드에서 입자 초기화 부분을 독립 함수로 분리.

\subsubsection{기능}

\begin{itemize}
  \item 양자상태 파일(\texttt{.txt}) 로드 및 Boltzmann 분포로 $(j, |m|)$ 샘플링
  \item 위치 초기화: 가우시안 분포 (x), 다이 레이저 공간 지터 (y, z)
  \item 속도 초기화: 빔 기울기 오프셋 + 가우시안 분포
  \item 레이저 지터 초기화: 공간/시간/위치 안정성 노이즈
\end{itemize}

\subsubsection{z-위치 초기화 공식}

\begin{lstlisting}
particle(:,3) = (30e-9 + laser(:,2)) .* z_init_vel_off ...
              - (60e-9 + laser(:,2)) .* particle(:,6) ...
              + laser(:,1);
% initial z from 30 ns delay + timing jitter
\end{lstlisting}

\subsection{\texttt{sim/time\_integrate.m} (신규)}

Velocity Verlet 시간 적분기. GPU/CPU 두 경로를 모두 지원.

\subsubsection{함수 시그니처}

\begin{lstlisting}
function velocity = time_integrate(position, velocity, Ij, Ips, ...
        I1_vec, I2_vec, aligned_cost, last_data, ...
        lambda, w0z, w0y, tau, MAX, tstep, t_ini, tend, use_gpu)
% use_gpu: true=GPU, false=CPU, omit=auto-detect
\end{lstlisting}

\subsubsection{CPU 경로 (parfor 내부용)}

\begin{lstlisting}
t = t_ini;
while t < tend
    position = position + velocity .* tstep;
    [~, last_data, a, ~] = axay_opt_partial( ...
        position, t + Ij, I1_vec, I2_vec, ...
        aligned_cost, last_data, Ips, lambda, w0z, w0y, tau, MAX);
    velocity = velocity + a .* tstep;
    t = t + tstep;
end
\end{lstlisting}

\subsubsection{GPU 경로 (단독 실행 시)}

\begin{lstlisting}
g_dev = gpuDevice;
bytes_per_particle = nI * 8 + 250;
gpu_batch_size = floor(g_dev.AvailableMemory * 0.8 / bytes_per_particle);
% auto-chunk by available VRAM
\end{lstlisting}

GPU 경로는 현재 parfor worker에서는 \texttt{false}로 비활성화하여 사용한다.
입자 수가 크게 늘어날 경우 재검토 가능.

\subsection{\texttt{sim/image\_processing.m} (신규)}

\subsubsection{설계}

\begin{itemize}
  \item 전자 반동(electron recoil) 속도 추가: 모든 그룹에 동일한 난수 적용
        (그룹 간 상관 노이즈, 실험 재현)
  \item PSF(\texttt{voigt2d\_aniso\_rot}) 1회 계산 후 모든 그룹에 재사용
  \item \texttt{hist3} + \texttt{conv2} 를 \texttt{parfor}로 병렬화
\end{itemize}

\begin{lstlisting}
parfor g = 1:n_groups
    h = hist3([vx_pix{g}, vy_pix{g}], 'Ctrs', {xgrid, ygrid});
    img_matrix{g} = conv2(h, M, 'same');
end
\end{lstlisting}

%% ---------------------------------------------------------------------------
\section{STN\_constants.mat 값 정리}

\begin{center}
\begin{tabular}{llp{6cm}}
\toprule
변수명 & 값 & 비고 \\
\midrule
\texttt{MAX}           & 20,000          & \texttt{run\_all.m}에서 200,000으로 오버라이드 \\
\texttt{I1\_list}      & 0.2             & 스칼라 \\
\texttt{I2\_list}      & [1.8, 3.2, 5.0, 7.2, 9.8, 12.8, 16.2, 20.0] & 1$\times$8 \\
\texttt{t\_ini}        & $-30\times10^{-9}$ s & \\
\texttt{tend}          & $30\times10^{-9}$ s & \\
\texttt{tstep}         & $10\times10^{-12}$ s & 총 6,000 스텝 \\
\texttt{jm\_data}      & \texttt{'alpha\_T1\_J=2n,j,m.txt'} & 파일명만 저장 \\
\midrule
\multicolumn{3}{l}{\textit{run\_all.m에 하드코딩된 이미지 처리 파라미터}} \\
\texttt{rep\_std}      & 4.0 m/s         & 전자 반동 std \\
\texttt{rep\_vel}      & 2.7 m/s         & 전자 반동 크기 \\
\texttt{vel\_pix}      & 2.28 m/s/pixel  & 속도-픽셀 변환 \\
\texttt{atten\_factorx}& $1/0.9278$      & PSF x축 \\
\texttt{atten\_factory}& $1/0.9453$      & PSF y축 \\
\bottomrule
\end{tabular}
\end{center}

%% ---------------------------------------------------------------------------
\section{알려진 한계 및 향후 개선 방안}

\subsection{현재 한계}

\begin{itemize}
  \item RAM 사용률 $\sim$85\% (64 GB 기준 $\sim$54 GB): batch\_size 증가 불가
  \item Worker 8개가 메모리 대역폭 포화점: 추가 병렬화 불가
  \item GPU는 현재 20,000 입자 규모에서 효과 없음
\end{itemize}

\subsection{향후 개선 가능 방안}

\begin{enumerate}[itemsep=4pt]
  \item \textbf{linearly\_2d + aligned\_cost를 single precision으로 변환}\\
        메모리 85\% $\to$ $\sim$55\%로 감소, 예상 속도 10--15\% 향상\\
        \textbf{주의}: position, velocity는 6,000 스텝 누적 오차로 인해 single 변환 불가

  \item \textbf{MAX를 200,000 이상으로 늘릴 경우}\\
        batch\_size / n\_workers 재조정 필요\\
        GPU 경로(\texttt{time\_integrate.m}의 \texttt{use\_gpu=true}) 재검토

  \item \textbf{axay\_opt\_partial.m 추가 최적화}\\
        \texttt{d\_alpha\_dI} 계산에서 \texttt{valid\_prev} 처리 분기를 제거하고
        첫 스텝에 last\_data를 명시적 초기화하는 방식으로 변경 가능
\end{enumerate}

%% ---------------------------------------------------------------------------
\section{빠른 참조}

\subsection{파라미터 튜닝 위치}

\begin{center}
\begin{tabular}{lll}
\toprule
파라미터 & 파일 & 줄 \\
\midrule
\texttt{n\_workers} & \texttt{run\_all.m} & 15 \\
\texttt{MAX} (분자 수) & \texttt{run\_all.m} & 36 \\
\texttt{batch\_size} & \texttt{run\_all.m} & 65 \\
이미지 처리 파라미터 & \texttt{run\_all.m} & 39--43 \\
\bottomrule
\end{tabular}
\end{center}

\subsection{실행 방법}

\begin{lstlisting}
% Run from MATLAB Command Window or Editor:
run('run_all.m')
% result saved to: result_alpha_T1_J=2n,j,m.mat
\end{lstlisting}

\end{document}
